\section{From a Gene to an Identifiable CAD}
\subsection{A Gene Is a Contract, Not Merely a Number}
A gene needs an identity, type, unit, legal domain, decoder relationship and version. A serialized number alone cannot distinguish a length from a dimensionless descriptor, discrete choice or activation coefficient. These semantics determine whether the decoded geometry is interpretable.

Jarambit's representation is defined in \src{crates/crescent/src/genome.rs}{the genome source}, with construction in \src{crates/crescent/src/pheno.rs}{the phenotype source}. It contains 128 numerical values and a material field. These are not 128 independent physical degrees of freedom: decoder dependencies, inactive features and derived geometry couple their effects. Material metadata has physical consequences and is distinct from an interpolated numeric coordinate.

\subsection{Decoder and Phenotype}
The decoder maps the stored representation to a planform and spatially varying thickness representation, then to CAD export and analysis paths. The intermediate phenotype connects parameter semantics to geometry. A decoder revision can change that geometry while the input CSV remains identical.

Candidate identity consequently includes the genome, material interpretation, decoder revision and export configuration. STEP, STL and analysis meshes have their own content identities. A shared filename does not establish equivalence. In particular, a render mesh and a mechanics mesh can differ; a candidate-specific numerical image requires the corresponding field/model binding.

\subsection{Export and Assessment Identity}
STEP and STL encode different representations with different topology checks. Unit interpretation, B-rep validity, triangle-shell closure and content association answer separate questions. Mass from raster/thickness approximation, mesh volume or a B-rep also depends on the representation and density interpretation. Similar numbers do not establish identical computation.

An assessment identifies the representation measured and the conversion context. Section~\ref{sec:fitness}'s export diagram illustrates these distinctions using saved records. Historical favorable flags remain observations under their original reader; they do not substitute for a current regeneration or manufacturing study.

\subsection{Designing a Gene: A Software Methodology}
A representation proposal states its purpose, typed meaning, serialization and decoder effect, then defines invalid-input handling and compatibility with older files. Observable software controls distinguish a correct implementation from a plausible output: a serialization round trip preserves semantics, a dimensional check detects incompatible interpretation, and a decoder-reference comparison detects a silently ignored parameter. Benign reference geometry can demonstrate those contracts without asserting application-specific performance.

Representation, objective and numerical-model versions remain separate. Otherwise an observed score change cannot be attributed to the representation. A migration records its mapping or marks incompatibility; it does not silently reinterpret historical observations. This methodology specifies software contracts rather than new physical parameters or shape modifications.

\subsection{What Complex Genomes Add}
Hierarchies, relationships, reusable substructures and spatial fields expand representational scope while creating hidden coupling and invalid combinations. Agents could assist in proposing or inspecting these relationships. Expressiveness would then be assessed through reliable decoding, interpretable variation and the questions the resulting phenotypes can answer, rather than parameter count alone.

The present evidence describes an existing authored representation and its traceable artifacts. Whether agent-proposed representations improve coverage or assessment remains a prospective research comparison; superior autonomous genome design is not demonstrated here.
